% Existing labels are retained to preserve all manuscript cross-references.
% The relocated material is highlighted in full.
\begingroup\color{blue}
\section{\rev{Task Relaxation with Soft and Hard Specifications}}
\label{ch02:sec:task_relaxation}

\begin{revision}
The preceding case study illustrates nominal single-robot synthesis.
When a task is infeasible, the first construction below allows controlled
specification violation. The second separates hard requirements from soft
preferences. Both constructions concern one robot. Their relaxed costs and
hard acceptance constraints supply the local planning models used by each
robot in Section~\ref{ch03:sec:multi-partial}, where new workspace information
can trigger plan revision.
\end{revision}

\subsection{\rev{Infeasible tasks}}\label{ch03:sec:infeasible}

\rev{Notation follows Sections~\ref{ch02:sec:system_abstraction_task_specification}
and~\ref{ch02:sec:solution}.} Consider a robot with local task specification~$\varphi$. By
Definition~\ref{ch02:def:feasible}, $\varphi$ may be infeasible over the
current abstraction~$\mathcal{T}_c$, either because the task is intrinsically
infeasible or because the model is incomplete. When the nominal synthesis
procedure of Section~\ref{ch02:sec:solution} fails, the objective becomes to
relax the specification and synthesize a plan that violates the original task as
little as possible.
A related approach by~\citet{kim2012approximate} constructs a relaxed
automaton~$\mathcal{A}_{\varphi}'$ that is close to
$\mathcal{A}_{\varphi}$ and feasible over~$\mathcal{T}_c$. The present
treatment uses a more discriminating objective: among accepting runs in the
relaxed product, select one that best preserves the original specification
under an explicit violation measure, and then infer the associated automaton
relaxation from that run.

\begin{problem}
Assume that the task specification is infeasible. How should the specification
be relaxed, and how can a discrete plan be synthesized so that it satisfies the
original task as much as possible?
\end{problem}

\subsubsection{Relaxed product automaton}
\label{ch03:subsec:relaxed-prod}

By Definition~\ref{ch02:def:feasible}, infeasibility of~$\varphi$ means that
the standard product automaton~$\mathcal{A}_p$ admits no accepting run. The
idea is therefore to relax the automaton-side constraints while keeping
the workspace unchanged. This leads to the relaxed product
B\"uchi automaton below.

\begin{definition}\label{ch03:def:relaxpro}
The relaxed product B\"uchi automaton is defined as:
\begin{equation}
\mathcal{A}_r \triangleq \mathcal{T}_c \times \mathcal{A}_{\varphi}
= (Q',\, 2^{AP},\, \delta',\, Q_0',\, \mathcal{F}',\, W_r),
\label{ch03:eq:relaxed-pba}
\end{equation}
where $Q' \triangleq \Pi \times Q$,
$AP \triangleq \{a_1,\, a_2,\, \cdots,\, a_K\}$,
$Q_0' \triangleq \Pi_0 \times Q_0$, and
$\mathcal{F}' \triangleq \Pi \times \mathcal{F}$. The transition relation
$\delta' : Q' \rightarrow 2^{Q'}$ satisfies
$\langle \pi_j,\, q_n \rangle \in \delta'(\langle \pi_i,\, q_m \rangle)$ if and
only if $(\pi_i,\, \pi_j) \in \longrightarrow_c$ and there exists
$l \in 2^{AP}$ such that $q_n \in \delta(q_m,\, l)$. The function
$W_r : Q' \times Q' \rightarrow \mathbb{R}^{+}$ is defined below.
\end{definition}

Relative to the standard product automaton in
Definition~\ref{ch02:def:product}, Definition~\ref{ch03:def:relaxpro} weakens
the automaton-side transition condition and augments the transition cost. The
weight now penalizes the distance between the observed label and the set of
labels that would enable the desired automaton transition.
\rev{To define this penalty, let
$\texttt{Eval} : 2^{AP} \rightarrow \{0,\, 1\}^{K}$ map each symbol to its
Boolean encoding under the fixed ordering $(a_1,\cdots,a_K)$, namely
$\texttt{Eval}(l)_k\triangleq 1$ if $a_k\in l$ and $0$ otherwise, for
$k\in\{1,\cdots,K\}$.} For $l,l' \in 2^{AP}$, define
$\rho(l,\, l') \triangleq \|\boldsymbol{\nu} - \boldsymbol{\nu}'\|_{1}$,
where $\boldsymbol{\nu} = \texttt{Eval}(l)$ and
$\boldsymbol{\nu}' = \texttt{Eval}(l')$. Thus $\rho$ is the Hamming distance
between the two encodings. For any nonempty set $\chi \subseteq 2^{AP}$, define
\begin{equation}
\texttt{Dist}(l,\, \chi) \triangleq
\begin{cases}
0, & \text{if } l \in \chi, \\
\min_{l' \in \chi} \rho(l,\, l'), & \text{otherwise},
\end{cases}
\label{ch03:eq:setdist}
\end{equation}
where $\chi$ is nonempty. The distance vanishes exactly when $l \in \chi$.
The transition weight of the relaxed product automaton is given by:
\begin{equation}
\begin{aligned}
W_r(\langle \pi_i,\, q_m \rangle,\,
\langle \pi_j,\, q_n \rangle)
\triangleq {} & w_c(\pi_i,\, \pi_j) \\
& {} + \alpha \cdot
\texttt{Dist}(L_c(\pi_i),\, \chi(q_m,\, q_n)),
\end{aligned}
\label{ch03:eq:WR}
\end{equation}
where $\langle \pi_j,\, q_n \rangle \in
\delta'(\langle \pi_i,\, q_m \rangle)$ and $\alpha \geq 0$. Here,
$\chi(q_m,\, q_n) \triangleq
\{\, l \in 2^{AP} \mid q_n \in \delta(q_m,\, l) \,\}$ collects all alphabets
that enable the transition from $q_m$ to $q_n$ in
$\mathcal{A}_{\varphi}$. By Definition~\ref{ch03:def:relaxpro}, this set is
nonempty whenever the relaxed transition exists. Thus
\eqref{ch03:eq:WR} combines implementation cost with a penalty for violating
the original automaton constraint; larger $\alpha$ gives stronger preference to
preserving the original task.

\begin{revision}
The relation $\longrightarrow_c$ determines which motions are feasible.
In \eqref{ch03:eq:WR}, $\alpha$ sets the motion cost assigned to one
proposition mismatch. Reducing the accumulated distance by one unit offsets
$\alpha$ units of accumulated motion cost. The resulting nonnegative scalar
edge weights allow the same prefix-suffix search as in the standard product.
A lexicographic objective minimizes violation first and breaks ties by motion
cost. A proposition-weighted distance assigns different costs to different
proposition mismatches. Section~\ref{ch03:sec:soft-hard} enforces hard task
requirements through both the transition guards and the B\"uchi acceptance
condition.
\end{revision}

\begin{example}\label{ch03:examp:dist}
Consider the NBA introduced in Example~\ref{ch02:examp:NBA} and the transition from $q_1$
to $q_3$. In this case,
$\chi(q_1,\, q_3)
= \{\{a_2\},\, \{a_1,\, a_2\},\, \{a_2,\, a_3\},\,
\{a_1,\, a_2,\, a_3\}\}$. Therefore,
$\textup{\texttt{Dist}}(\{a_1\},\, \chi(q_1,\, q_3))
= \rho(\{a_1\},\, \{a_1,\, a_2\}) = 1$,
$\textup{\texttt{Dist}}(\{a_2\},\, \chi(q_1,\, q_3))
= \rho(\{a_2\},\, \{a_2\}) = 0$, and
$\textup{\texttt{Dist}}(\{a_2,\, a_3\},\, \chi(q_1,\, q_3))
= \rho(\{a_2,\, a_3\},\, \{a_2,\, a_3\}) = 0$.
\end{example}

\subsubsection{Balanced accepting run}\label{ch03:subsec:probStatement}

\rev{The relaxed product $\mathcal{A}_r$ can be searched regardless of whether
the standard product $\mathcal{A}_p$ admits an accepting run.}
\begin{revision}
Different accepting runs may have different motion costs and specification
violations, so the prefix-suffix structure in
Section~\ref{ch02:subsec:optimality} is retained:
\[
R \triangleq q_0' q_1' \cdots q_{f-1}'
[\, q_f' q_{f+1}' \cdots q_n' \,]^{\omega}
= \langle \pi_0,\, q_0 \rangle \cdots
\langle \pi_{f-1},\, q_{f-1} \rangle
[\, \langle \pi_f,\, q_f \rangle
\langle \pi_{f+1},\, q_{f+1} \rangle \cdots
\langle \pi_n,\, q_n \rangle \,]^{\omega},
\]
where $q_0' = \langle \pi_0,\, q_0 \rangle \in Q_0'$ and
$q_f' = \langle \pi_f,\, q_f \rangle \in \mathcal{F}'$.
\end{revision}
The total cost of~$R$ in the relaxed product automaton is defined as:
\begin{equation}
\texttt{Cost}(R,\, \mathcal{A}_r)
\triangleq \texttt{cost}_{\tau}
+ \alpha \cdot \texttt{dist}_{\varphi},
\label{ch03:eq:prod-cost2}
\end{equation}
where $\tau \triangleq R|_{\Pi}$ is the induced motion plan. \rev{The implementation
cost in \eqref{ch03:eq:prod-cost2} is given by}
\begin{equation}
\color{blue}
\texttt{cost}_{\tau} \triangleq
\sum_{i=0}^{f-1}w_c(\pi_i,\pi_{i+1})
+\gamma\left(\sum_{i=f}^{n-1}w_c(\pi_i,\pi_{i+1})
+w_c(\pi_n,\pi_f)\right),
\label{ch03:eq:cost}
\end{equation}
\rev{where the first sum covers the prefix and the $\gamma$-weighted term
covers one complete suffix cycle. The accumulated specification distance is}
\begin{equation}
\color{blue}
\begin{aligned}
\texttt{dist}_{\varphi}\triangleq {}&
\sum_{i=0}^{f-1}\texttt{Dist}(L_c(\pi_i),\chi(q_i,q_{i+1})) \\
&+\gamma\sum_{i=f}^{n-1}\texttt{Dist}(L_c(\pi_i),\chi(q_i,q_{i+1})) \\
&+\gamma\,\texttt{Dist}(L_c(\pi_n),\chi(q_n,q_f)),
\end{aligned}
\label{ch03:eq:dist}
\end{equation}
\begin{revision}
where $\gamma>0$ weights the relative importance of the transient prefix
and the recurrent suffix. The closing edge of the suffix is included,
as in \eqref{ch02:eq:prod-cost}. Positivity of $\gamma$ is required for
zero accumulated distance to imply that every recurrent transition
satisfies the original automaton constraint.
\end{revision}

\begin{problem}\label{ch03:pro:balanced}
Find an accepting run of $\mathcal{A}_r$ that minimizes the cost in
\eqref{ch03:eq:prod-cost2}.
\end{problem}

The solution of Problem~\ref{ch03:pro:balanced} is the balanced accepting run
$R_{\textup{bal}}$ of~$\mathcal{A}_r$, with induced plan
$\tau_{\textup{bal}} \triangleq R_{\textup{bal}}|_{\Pi}$.
\rev{If the guard for $\chi$ is represented by an ordered binary decision
diagram, \eqref{ch03:eq:setdist} is the minimum cost of a path from the root
to the true terminal. Each branch costs one if its bit differs from the
corresponding bit of $\texttt{Eval}(l)$ and zero otherwise. Propositions skipped
on the path can retain their bits in $\texttt{Eval}(l)$.}


\subsubsection{Balanced plan synthesis}\label{ch03:subsec:singleFeedback}

For fixed $\alpha$ and $\gamma$, the relaxed automaton $\mathcal{A}_r$ may be
built explicitly from Definition~\ref{ch03:def:relaxpro} or generated on the
fly. Algorithm~\ref{ch02:alg:optimal} then computes
$R_{\textup{bal}}$ and the plan
$\tau_{\textup{bal}} = R_{\textup{bal}}|_{\Pi}$. Although
$\mathcal{A}_r$ admits more automaton transitions than $\mathcal{A}_p$, it does
not relax the transition relation of~$\mathcal{T}_c$; hence every balanced plan
is implementable.
Once $R_{\textup{bal}}$ is known, $\texttt{cost}_{\tau}$ and
$\texttt{dist}_{\varphi}$ follow from \eqref{ch03:eq:cost} and
\eqref{ch03:eq:dist}. A relaxed automaton~$\mathcal{A}_{\varphi}'$ can also be
obtained by adding, along $R_{\textup{bal}}$, the automaton edges needed to
support the observed labels, yielding a relaxation in the sense of
\citet{kim2012approximate}.
The parameter~$\alpha$ controls the trade-off between implementation cost and
specification violation. In practice, Algorithm~\ref{ch02:alg:optimal} can be
run for several values of $\alpha$, after which distinct balanced runs are
compared through $(\texttt{cost}_{\tau},\, \texttt{dist}_{\varphi})$ and the
selected plan is executed by Algorithm~\ref{ch02:alg:hybrid}.

\begin{lemma}\label{ch03:lem:distanceProof}
If $\textup{\texttt{dist}}_{\varphi} = 0$, then
$\tau_{\textup{bal}} \models \varphi$.
\end{lemma}

\begin{proof}
By \eqref{ch03:eq:setdist}, the function $\texttt{Dist}(\cdot,\, \cdot)$ is
nonnegative. Therefore, $\texttt{dist}_{\varphi} = 0$ implies that
$q_n \in \delta(q_m,\, L_c(\pi_i))$ for every transition
$(\langle \pi_i,\, q_m \rangle,\,
\langle \pi_j,\, q_n \rangle)$ along $R_{\textup{bal}}$. Hence every
transition of $R_{\textup{bal}}$ also satisfies the transition condition of the
standard product automaton $\mathcal{A}_p$ from
Definition~\ref{ch02:def:product}. Since $\mathcal{A}_p$ and $\mathcal{A}_r$
share the same state set, the same initial states, and the same accepting
states, $R_{\textup{bal}}$ is also an accepting run of~$\mathcal{A}_p$.
Lemma~\ref{ch02:lem:satisfy} then yields
$\tau_{\textup{bal}} \models \varphi$.
\end{proof}


\begin{revision}
\begin{theorem}\label{ch03:thm:balanced-feasible}
Suppose that $\varphi$ is feasible over $\mathcal{T}_c$ and $\gamma>0$.
Let $R_{\textup{feas}}$ be any accepting prefix-suffix run of the standard
product $\mathcal{A}_p$. Evaluate its motion cost
$C_{\textup{feas}}\triangleq J(R_{\textup{feas}})$ using
\eqref{ch02:eq:prod-cost} with the same suffix weight $\gamma$ as in the
relaxed objective, and define
\begin{equation}
\underline{\alpha}\triangleq
\frac{C_{\textup{feas}}}{\min\{1,\gamma\}},
\label{ch03:underalpha}
\end{equation}
where $\min\{1,\gamma\}$ is a lower bound on any positive accumulated distance.
If $\alpha>\underline{\alpha}$, every minimizing relaxed run has zero
accumulated distance, and its projected plan satisfies $\varphi$.
\end{theorem}

\begin{proof}
The run $R_{\textup{feas}}$ is also accepting in $\mathcal{A}_r$ and has
zero accumulated distance, so its relaxed cost is $C_{\textup{feas}}$.
Every transition distance is a nonnegative integer. A positive term in
\eqref{ch03:eq:dist} contributes at least $1$ on the prefix or $\gamma$
on the suffix. Since motion costs are nonnegative, any run with positive
distance costs at least $\alpha\min\{1,\gamma\}$. For
$\alpha>\underline{\alpha}$ this lower bound exceeds $C_{\textup{feas}}$.
Every minimizing run therefore has zero accumulated distance, and its
projection satisfies $\varphi$ by Lemma~\ref{ch03:lem:distanceProof}.
\end{proof}

To construct $R_{\textup{feas}}$, find a path from an initial product state
to an accepting state in $\mathcal{A}_p$ and a nonempty cycle through it.
Compute $C_{\textup{feas}}$ including the cycle's closing edge and choose
$\alpha$ strictly above the value in \eqref{ch03:underalpha}. At equality, a run with
positive distance may tie with the feasible run. A lower-cost feasible run
gives a tighter sufficient bound. If an exhaustive search finds no reachable
accepting cycle, the task is infeasible. A time-limited search that returns
no such run leaves feasibility unresolved.
\end{revision}



\begin{revision}
\begin{example}\label{ch03:examp:singleRS}
Consider the four-region graph in Fig.~3 of \citet{guo2012infeasible}.
Its routes from $\pi_0$ to $\pi_3$ go through $\pi_1$ with edge costs $10,30$
or through $\pi_2$ with edge costs $20,40$. The self-loops at $\pi_0$ and
$\pi_3$ each cost $5$. The labels of $\pi_0,\pi_1,\pi_2,\pi_3$ are
$\emptyset$, $\{a_2,a_3\}$, $\{a_2\}$, and $\{a_1\}$, respectively.

The task is $\varphi=\Diamond\Box a_1\land\Box(\neg a_2\land\neg a_3)$.
In a two-state automaton for this task, the initial state $q_0$ has a
self-loop guarded by $\neg a_2\land\neg a_3$. Both transitions into the
accepting state $q_1$ require $a_1\land\neg a_2\land\neg a_3$.

With $\gamma=5$, the balanced plans for increasing $\alpha$ are
$[\pi_0]^\omega$, $\pi_0\pi_1[\pi_3]^\omega$, and $\pi_0\pi_2[\pi_3]^\omega$.
Their motion costs and accumulated distances are $(30,6)$, $(70,2)$,
and $(90,1)$, respectively. For the last two plans, the transition from
$q_0$ to $q_1$ reads $\{a_1\}$ on the first self-loop at $\pi_3$, which costs
$5$ on the prefix. The recurrent self-loop contributes $5\gamma$ to each
motion cost. The plans exchange optimality at $\alpha=10$ and $\alpha=20$.
\end{example}
\end{revision}












%%%%%%%%%%%%%%%%%%%%%%%%%%%%%%%%%%%%%%%%%%%%%%%%%%%
\subsection{\rev{Soft and hard specifications}}\label{ch03:sec:soft-hard}

The previous subsection allowed the entire specification to be relaxed. Many
applications instead separate safety-critical requirements, which must never be
violated, from performance or achievement objectives, which may be relaxed under
partial workspace knowledge. This subsection develops a synthesis framework for
that hard--soft structure.

\subsubsection{Problem formulation}\label{ch03:subsec:formulation}

The robot transition system is still denoted by $\mathcal{T}_c$, as in
Definition~\ref{ch02:def:FTS}. The task specification remains an LTL formula
over $AP$, but now has the form
\begin{equation}
\varphi \triangleq \varphi^{\textup{soft}} \wedge \varphi^{\textup{hard}},
\label{ch03:eq:spec}
\end{equation}
where $\varphi^{\textup{hard}}$ and $\varphi^{\textup{soft}}$ are the hard and
soft components. Hard requirements include constraints such as obstacle
avoidance or recurrent charging; soft requirements encode objectives whose
feasibility may depend on incomplete workspace information. Thus relaxation is
allowed only for the soft component.

\begin{problem}
Given the specification in \eqref{ch03:eq:spec}, synthesize a motion and task
plan such that $\varphi^{\textup{hard}}$ is satisfied fully, while
$\varphi^{\textup{soft}}$ is satisfied as much as possible.
\end{problem}

\subsubsection{Safety-ensured product automaton}
\label{ch03:subsec:intersection}

The decomposition in \eqref{ch03:eq:spec} leads naturally to three notions of
correctness.

\begin{definition}\label{ch03:def:vs}
A finite or infinite path $\tau$ of $\mathcal{T}_c$ is called \emph{valid}
if every consecutive pair of states along $\tau$ belongs to
$\longrightarrow_c$. An infinite valid path $\tau$ is called \emph{safe} if
$\tau \models \varphi^{\textup{hard}}$, and \emph{satisfying} if
$\tau \models \varphi$.
\end{definition}

Let
$\mathcal{A}^{\textup{hard}}
= (Q_1,\, 2^{AP},\, \delta_1,\, Q_{1,0},\, \mathcal{F}_1)$ and
$\mathcal{A}^{\textup{soft}}
= (Q_2,\, 2^{AP},\, \delta_2,\, Q_{2,0},\, \mathcal{F}_2)$ be the NBAs for
$\varphi^{\textup{hard}}$ and $\varphi^{\textup{soft}}$, respectively. Define
$\chi_1(q_1,\, \check{q}_1) \triangleq
\{\, l \in 2^{AP} \mid \check{q}_1 \in \delta_1(q_1,\, l) \,\}$ and
$\chi_2(q_2,\, \check{q}_2) \triangleq
\{\, l \in 2^{AP} \mid \check{q}_2 \in \delta_2(q_2,\, l) \,\}$ for the hard
and soft automata. The hard component is enforced exactly, while the soft
component is relaxed in the following intersection automaton.

\begin{definition}\label{ch03:def:conj}
The relaxed intersection of $\mathcal{A}^{\textup{hard}}$ and
$\mathcal{A}^{\textup{soft}}$ is given by:
\begin{equation}
\tilde{\mathcal{A}}_{\varphi}
\triangleq
(Q,\, 2^{AP},\, \delta,\, Q_0,\, \mathcal{F}),
\label{ch03:eq:intersect}
\end{equation}
where
$Q \triangleq Q_1 \times Q_2 \times \{1,\, 2\}$,
$Q_0 \triangleq Q_{1,0} \times Q_{2,0} \times \{1\}$, and
$\mathcal{F} \triangleq \mathcal{F}_1 \times Q_2 \times \{1\}$. The transition
relation $\delta : Q \times 2^{AP} \rightarrow 2^Q$ satisfies
$\langle \check{q}_1,\, \check{q}_2,\, \check{t} \rangle \in
\delta(\langle q_1,\, q_2,\, t \rangle,\, l)$ if and only if
$l \in \chi_1(q_1,\, \check{q}_1)$,
$\chi_2(q_2,\, \check{q}_2) \neq \emptyset$, and either
$q_t \notin \mathcal{F}_t$ with $\check{t} = t$, or
$q_t \in \mathcal{F}_t$ with
$\check{t} = \textup{\texttt{mod}}(t,\, 2) + 1$.
\end{definition}

Definition~\ref{ch03:def:conj} relaxes only the soft component: the actual
alphabet must satisfy the hard transition, whereas the soft transition must
only be enabled by some alphabet. The auxiliary variable $t \in \{1,\, 2\}$
alternates the acceptance obligation as in the standard NBA intersection
\citep{baier2008principles}. Denote by $r|_{Q_1}$ and $r|_{Q_2}$ the projections
of a run~$r$ onto the hard and soft automata.

\begin{algorithm}[t]
\DontPrintSemicolon
\caption{Construction of the relaxed intersection, \texttt{RelaxInt}()}
\label{ch03:alg:relaxedint}
\KwIn{$\mathcal{A}^{\textup{soft}}$, $\mathcal{A}^{\textup{hard}}$}
\KwOut{$\tilde{\mathcal{A}}_{\varphi}$}

\tcc{\rev{Form states and alternate the hard and soft acceptance obligations.}}
\ForEach{$q_1 \in Q_1$, $q_2 \in Q_2$, $t \in \{1,\,2\}$}{
  $q_m = \langle q_1,\, q_2,\, t \rangle \in Q$\;

  \If{$q_1 \in Q_{1,0}$, $q_2 \in Q_{2,0}$, and $t = 1$}{
    $q_m \in Q_0$\;
  }

  \If{$q_1 \in \mathcal{F}_1$ and $t = 1$}{
    $q_m \in \mathcal{F}$\;
  }

  \ForEach{$\check{q}_1 \in \textup{\texttt{Post}}(q_1)$,
           $\check{q}_2 \in \textup{\texttt{Post}}(q_2)$,
           $\check{t} \in \{1,\,2\}$}{
    $q_n = \langle \check{q}_1,\, \check{q}_2,\, \check{t} \rangle \in Q$\;

    \If{$(q_1 \notin \mathcal{F}_1$ and $t = 1$ and $\check{t} = 1)$ or
        $(q_2 \notin \mathcal{F}_2$ and $t = 2$ and $\check{t} = 2)$ or
        $(q_1 \in \mathcal{F}_1$ and $t = 1$ and $\check{t} = 2)$ or
        $(q_2 \in \mathcal{F}_2$ and $t = 2$ and $\check{t} = 1)$}{
      $q_n \in \delta(q_m,\, l)$, for all $l \in \chi_1(q_1,\, \check{q}_1)$\;
    }
  }
}
\textbf{return} $\tilde{\mathcal{A}}_{\varphi}$
\end{algorithm}

\begin{theorem}\label{ch03:theo:saferun}
Given an accepting run $r$ of $\tilde{\mathcal{A}}_{\varphi}$, the projection
$r|_{Q_1}$ is an accepting run of $\mathcal{A}^{\textup{hard}}$. Moreover,
$\mathcal{L}_{\omega}(\tilde{\mathcal{A}}_{\varphi}) \subseteq
\mathcal{L}_{\omega}(\mathcal{A}^{\textup{hard}})$.
\end{theorem}

\begin{proof}
Since $r$ is accepting, it visits the set
$\mathcal{F} = \mathcal{F}_1 \times Q_2 \times \{1\}$ infinitely often.
Therefore, the projection $r|_{Q_1}$ visits $\mathcal{F}_1$ infinitely often.
In addition, every transition of $r$ satisfies
$l \in \chi_1(q_1,\, \check{q}_1)$ by
Definition~\ref{ch03:def:conj}. Hence every projected transition of $r|_{Q_1}$
is a valid transition of $\mathcal{A}^{\textup{hard}}$. It follows that
$r|_{Q_1}$ is an accepting run of $\mathcal{A}^{\textup{hard}}$.
For the language inclusion, let
$\sigma \in \mathcal{L}_{\omega}(\tilde{\mathcal{A}}_{\varphi})$. Then
$\sigma$ induces an accepting run $r_{\sigma}$ of
$\tilde{\mathcal{A}}_{\varphi}$. By the first part of the proof,
$r_{\sigma}|_{Q_1}$ is an accepting run of
$\mathcal{A}^{\textup{hard}}$. Therefore,
$\sigma \in \mathcal{L}_{\omega}(\mathcal{A}^{\textup{hard}})$, which proves
the desired inclusion.
\end{proof}

To preserve the hard specification while still optimizing the soft one, combine
$\tilde{\mathcal{A}}_{\varphi}$ with the workspace transition system in a
relaxed product automaton.

\begin{definition}\label{ch03:def:safe}
The safety-ensured relaxed product B\"uchi automaton is defined as:
\begin{equation}
\tilde{\mathcal{A}}_{r}
\triangleq
\mathcal{T}_c \times \tilde{\mathcal{A}}_{\varphi}
=
(Q',\, 2^{AP},\, \delta',\, Q_0',\, \mathcal{F}',\, W_r),
\label{ch03:eq:safe-product}
\end{equation}
where $Q' \triangleq \Pi \times Q$,
$Q_0' \triangleq \Pi_0 \times Q_0$, and
$\mathcal{F}' \triangleq \Pi \times \mathcal{F}$. The transition relation
$\delta' : Q' \rightarrow 2^{Q'}$ satisfies
$\langle \pi_j,\, q_n \rangle \in \delta'(\langle \pi_i,\, q_m \rangle)$ if and
only if $(\pi_i,\, \pi_j) \in \longrightarrow_c$ and
$q_n \in \delta(q_m,\, L_c(\pi_i))$. The weight function
$W_r : Q' \times Q' \rightarrow \mathbb{R}^{+}$ is given by:
\begin{equation}
W_r(\langle \pi_i,\, q_m \rangle,\,
\langle \pi_j,\, q_n \rangle)
\triangleq  w_c(\pi_i,\, \pi_j) + \alpha \cdot
\textup{\texttt{Dist}}(L_c(\pi_i),\, \chi_2(q_2,\, \check{q}_2)),
\label{ch03:eq:weight-soft}
\end{equation}
where $q_m = \langle q_1,\, q_2,\, t \rangle$,
$q_n = \langle \check{q}_1,\, \check{q}_2,\, \check{t} \rangle$, and
$\alpha \geq 0$.
\end{definition}

The weight in \eqref{ch03:eq:weight-soft} combines implementation cost with
the violation distance of the soft specification.
\rev{The transition relation of $\tilde{\mathcal{A}}_{\varphi}$ enforces
the hard guards, and its B\"uchi acceptance condition requires infinitely
many visits to hard accepting states. Together they ensure satisfaction of
the hard specification.}

\begin{theorem}\label{ch03:theo:safety}
Assume that $R$ is an accepting run of $\tilde{\mathcal{A}}_{r}$. Then its
projection $\tau \triangleq R|_{\Pi}$ is both valid and safe in the sense of
Definition~\ref{ch03:def:vs}.
\end{theorem}

\begin{proof}
The path $\tau$ is valid because every transition in $\delta'$ projects onto a
transition in $\longrightarrow_c$. Since $R$ is an accepting run of
$\tilde{\mathcal{A}}_{r}
= \mathcal{T}_c \times \tilde{\mathcal{A}}_{\varphi}$, the trace of $\tau$ is
accepted by $\tilde{\mathcal{A}}_{\varphi}$. By
Theorem~\ref{ch03:theo:saferun},
$\texttt{trace}(\tau) \in
\mathcal{L}_{\omega}(\mathcal{A}^{\textup{hard}})$. Then, by
Lemma~\ref{ch02:lem:LTL-BA},
$\texttt{trace}(\tau) \in \texttt{Words}(\varphi^{\textup{hard}})$, which is
equivalent to $\tau \models \varphi^{\textup{hard}}$. Hence, the resulting plan $\tau$ is safe.
\end{proof}

\rev{As in Section~\ref{ch03:sec:infeasible}, accepting runs use the
prefix-suffix form of \eqref{ch02:eq:prefix-suffix} with $\gamma>0$.}
Their total cost is
defined by:
\begin{equation}
\color{blue}
\begin{aligned}
\texttt{Cost}(R,\, \tilde{\mathcal{A}}_{r})
&\triangleq
\texttt{cost}_{\tau}
+ \alpha \cdot \texttt{dist}_{\varphi^{\textup{soft}}}, \\
\texttt{dist}_{\varphi^{\textup{soft}}}
&\triangleq
\sum_{i=0}^{f-1}\texttt{Dist}(L_c(\pi_i),
\chi_2(q_i'|_{Q_2},q_{i+1}'|_{Q_2})) \\
&\quad+\gamma\left(\sum_{i=f}^{n-1}\texttt{Dist}(L_c(\pi_i),
\chi_2(q_i'|_{Q_2},q_{i+1}'|_{Q_2}))
+\texttt{Dist}(L_c(\pi_n),\chi_2(q_n'|_{Q_2},q_f'|_{Q_2}))\right),
\end{aligned}
\label{ch03:eq:prod-cost3}
\end{equation}
where $\tau \triangleq R|_{\Pi}$, $\texttt{cost}_{\tau}$ is given by
\eqref{ch03:eq:cost}, and $q_i'|_{Q_2}$ denotes the projection of the product
state $q_i'$ onto the soft-automaton state set $Q_2$.

\begin{problem}\label{ch03:pro:safe}
Find an accepting run of $\tilde{\mathcal{A}}_{r}$ that minimizes the cost in
\eqref{ch03:eq:prod-cost3}.
\end{problem}

\subsubsection{Safe plan synthesis}\label{ch03:subsec:safe-plan-synthesis}

The solution of Problem~\ref{ch03:pro:safe} is the safe accepting run
$R_{\textup{safe}}$ of $\tilde{\mathcal{A}}_{r}$, with induced plan
$\tau_{\textup{safe}} \triangleq R_{\textup{safe}}|_{\Pi}$. For fixed
$\alpha$ and $\gamma$, $\tilde{\mathcal{A}}_{r}$ can be built explicitly or on
the fly. Algorithm~\ref{ch02:alg:optimal} then computes
$R_{\textup{safe}}$, and Theorem~\ref{ch03:theo:safety} guarantees that
$\tau_{\textup{safe}}$ is valid and safe. As before, varying $\alpha$ yields
candidate runs that can be compared through
$(\texttt{cost}_{\tau},\, \texttt{dist}_{\varphi^{\textup{soft}}})$ before
implementation by Algorithm~\ref{ch02:alg:hybrid}.

\begin{lemma}\label{ch03:lem:distanceProofsafe}
If $\textup{\texttt{dist}}_{\varphi^{\textup{soft}}} = 0$, then
$\tau_{\textup{safe}} \models \varphi$.
\end{lemma}

\begin{revision}
\begin{proof}
Since $\gamma>0$, zero accumulated soft distance implies that every
transition in $R_{\textup{safe}}$ satisfies its soft guard. The run
also satisfies every hard guard and the B\"uchi acceptance condition
of the hard--soft intersection. Restoring the soft guards preserves
each edge of this run as well as its initial and accepting states.
Thus $R_{\textup{safe}}$ is an accepting run of the unrelaxed product,
and Lemma~\ref{ch02:lem:satisfy} gives $\tau_{\textup{safe}}\models\varphi$.
\end{proof}
\end{revision}














%%%%%%%%%%%%%%%%%%%%%%%%%%%%%%%%%%%%%%%%%%%%%%%%%%%%
%%%%%%%%%%%%%%%%%%%%%%%%%%%%%%%%%%%%%%%%%%%%%%%%%%%
%%%%%%%%%%%%%%%%%%%%%%%%%%%%%%%%%%%%%%%%%%%%%%%%%%%%
\endgroup
